\section{Teorema Terletak di Atas dan Naik}
Sekarang kita menafsirkan keintegralan melalui geometri $\spec$.
Secara umum, untuk suatu homomorfisme gelanggang $R \to S$, peta terinduksi $\spec S \to
\spec R$ tidak harus berperilaku baik secara topologis; misalnya, meskipun $S$ merupakan
aljabar-$R$ terbangkitkan hingga, citra $\spec S$ di dalam $\spec R$ tidak harus
terbuka ataupun tertutup.\footnote{Namun, benar bahwa jika $R$
\emph{noetherian} (lihat \rref{noetherian}) dan $S$ terbangkitkan hingga di atas
$R$, maka citra $\spec S$ bersifat \emph{konstruktibel,} yakni gabungan hingga
himpunan-himpunan tertutup lokal. Ini adalah teorema konstruktibilitas Chevalley; buktinya berada di luar suplemen pilihan ini.}

Kita akan melihat bahwa lebih banyak hal dapat dikatakan di bawah syarat keintegralan.



\subsection{Teorema terletak di atas}

Secara umum, untuk suatu morfisme varietas aljabar $f: X \to Y$, citra
himpunan bagian tertutup $Z \subset X$ sama sekali belum tentu tertutup. Sebagai contoh, fungsi reguler $f: X \to \mathbb{C}$ 
yang merupakan pemetaan tertutup harus surjektif atau konstan (misalnya jika $X$
terhubung).
Meskipun demikian, kita dapat mengatakan lebih banyak di bawah hipotesis keintegralan.

\begin{proposition}[Teorema terletak di atas] \label{lyingover}
Jika $\phi: R \to R'$ merupakan morfisme integral, maka peta terinduksi
\[  \spec R' \to \spec R  \]
merupakan pemetaan tertutup; peta ini surjektif jika $\phi$ injektif.
 \end{proposition}

Cara lain untuk menyatakan klaim terakhir, tanpa menyebut $\spec R'$, adalah sebagai
berikut. Andaikan $\phi$ injektif dan integral. Jika
$\mathfrak{p} \subset R$ merupakan ideal prima, maka terdapat $\mathfrak{q} \subset R'$
sedemikian sehingga $\mathfrak{p}$ adalah prapeta $\phi^{-1}(\mathfrak{q})$.

\begin{proof} Pertama andaikan $\phi$ injektif; dalam hal ini kita harus membuktikan bahwa
peta $\spec R' \to \spec R$ surjektif.
Mari kita reduksi pembuktian ke kasus gelanggang lokal.
Untuk suatu ideal prima $\mathfrak{p} \in \spec R$, kita harus menunjukkan bahwa $\mathfrak{p}$
muncul sebagai prapeta suatu elemen $\spec R'$.
Jadi, kita ganti $R$ dengan $R_{\mathfrak{p}}$. Kita memperoleh pemetaan
\[ \phi_{\mathfrak{p}}: R_{\mathfrak{p}} \to (R- \mathfrak{p})^{-1} R'  \]
yang injektif jika $\phi$ injektif, karena pelokalan merupakan funktor eksak. Di sini
kita
telah melokalkan $R, R'$ pada himpunan multiplikatif $R - \mathfrak{p}$.

Perhatikan bahwa $\phi_{\mathfrak{p}}$ juga merupakan perluasan integral. Hal ini mengikuti
karena keintegralan dipertahankan oleh perubahan basis.
Sekarang kita akan membuktikan hasil tersebut untuk $\phi_{\mathfrak{p}}$; khususnya, kita
akan menunjukkan
bahwa terdapat ideal prima dari $(R- \mathfrak{p})^{-1} R'$ yang prapetanya adalah
$\mathfrak{p}R_{\mathfrak{p}}$. Hal ini akan menyiratkan bahwa apabila ideal prima tersebut
ditarik kembali ke $R'$, prapetanya adalah $\mathfrak{p}$ di dalam $R$. Secara terperinci, kita
dapat mempertimbangkan diagram
\[ \xymatrix{
\spec (R-\mathfrak{p})^{-1} R'\ar[d]  \ar[r] & \spec R_{\mathfrak{p}} \ar[d] \\
\spec R' \ar[r] &  \spec R
}\]
yang menunjukkan bahwa jika $\mathfrak{p} R_{\mathfrak{p}}$ muncul dalam citra pemetaan atas, maka
$\mathfrak{p}$
muncul sebagai citra sesuatu di $\spec R'$.
Jadi, untuk membuktikan proposisi (yakni kasus ketika $\phi$
injektif), cukup menangani kasus $R$ lokal dan
$\mathfrak{p}$ merupakan ideal maksimal.


Dengan kata lain, kita perlu menunjukkan bahwa:
\begin{quote}
Jika $R$ merupakan gelanggang \emph{lokal} dan $\phi: R \hookrightarrow R'$ suatu
morfisme integral injektif, maka ideal maksimal $R$ merupakan prapeta
suatu elemen $\spec R'$.
\end{quote}

Andaikan $R$ lokal dengan ideal maksimal $\mathfrak{p}$. Kita ingin mencari ideal prima $\mathfrak{q} \subset R'$ sedemikian sehingga
$\mathfrak{p}  = \phi^{-1}(\mathfrak{q})$. Karena $\mathfrak{p}$ sudah
maksimal, cukup ditunjukkan bahwa $\mathfrak{p} \subset
\phi^{-1}(\mathfrak{q})$. Secara khusus, kita perlu menunjukkan bahwa terdapat ideal
prima $\mathfrak{q}$ sedemikian sehingga
\( \mathfrak{p} R' \subset \mathfrak{q}.  \)
Prapeta ideal ini adalah $\mathfrak{p}$.

Jika $\mathfrak{p}R' \neq R'$, maka
$\mathfrak{q}$ ada, karena setiap ideal sejati suatu gelanggang termuat dalam
ideal maksimal. Bahkan kita akan menunjukkan
\begin{equation} \label{thingdoesn'tgenerate} \mathfrak{p} R' \neq R',
\end{equation}
atau bahwa $\mathfrak{p}$ tidak membangkitkan ideal satuan dalam $R'$.
Jika kita membuktikan \eqref{thingdoesn'tgenerate}, kita dapat menemukan
$\mathfrak{q}$ yang dicari, dan pembuktian selesai.

Andaikan
sebaliknya, yakni $\mathfrak{p}R' = R'$. Kita akan memperoleh kontradiksi menggunakan
Lema Nakayama (\cref{nakayama}). 
Saat ini kita tidak dapat menerapkan Lema Nakayama secara langsung karena $R'$ bukan
modul-$R$ terbangkitkan hingga. Gagasannya ialah “menurunkan” “bukti” bahwa 
\eqref{thingdoesn'tgenerate} gagal ke suatu subaljabar kecil dari $R'$, lalu
memperoleh kontradiksi.
Untuk melakukannya, perhatikan bahwa $1 \in \mathfrak{p}R'$, dan kita dapat menulis
\[ 1 = \sum x_i \phi(y_i)  \]
dengan $x_i \in R', y_i \in \mathfrak{p}$.
Inilah “bukti” bahwa \eqref{thingdoesn'tgenerate} gagal, dan bukti tersebut
hanya melibatkan data dalam jumlah hingga. 

Misalkan $R''$ subaljabar dari $R'$ yang dibangkitkan oleh $\phi(R)$ dan semua $x_i$. Maka
$R'' \subset R'$ dan merupakan aljabar-$R$ terbangkitkan hingga, karena dibangkitkan oleh
semua $x_i$. Akan tetapi, $R''$ integral di atas $R$, sehingga terbangkitkan hingga sebagai
modul-$R$ berdasarkan \cref{finitelygeneratedintegral}. Di sinilah
keintegralan berperan.

Jadi, $R''$ merupakan modul-$R$ terbangkitkan hingga. Selain itu, persamaan
$1 = \sum x_i \phi(y_i)$ menunjukkan bahwa $\mathfrak{p}R'' = R''$. Namun, hal ini
bertentangan dengan Lema Nakayama. Kontradiksi tersebut menunjukkan bahwa
$\mathfrak{p}$ tidak dapat membangkitkan $(1)$ dalam $R'$, dan membuktikan bagian surjektivitas
dari teorema terletak di atas.

Terakhir, kita perlu menunjukkan bahwa jika $\phi: R \to R'$ adalah \emph{sebarang} morfisme
integral, maka $\spec R' \to \spec R$ merupakan pemetaan tertutup. Misalkan $X = V(I) $ suatu
himpunan bagian tertutup dari $\spec R'$. Maka citra $X$ di dalam $\spec R$ adalah citra
pemetaan
\[ \spec R'/I \to \spec R   \]
yang diperoleh dari morfisme $R \to R' \to  R'/I$, yang integral; dengan demikian, kita
direduksi untuk menunjukkan bahwa setiap morfisme integral $\phi$ mempunyai citra tertutup pada
$\spec$. 
Jadi, kita direduksi ke $X = \spec R'$ jika $R'$ dikesampingkan dan diganti dengan
$R'/I$.

Dengan kata lain, kita harus membuktikan pernyataan berikut. Misalkan $\phi: R \to R'$ suatu
morfisme integral; maka citra $\spec R'$ di dalam $\spec R$ tertutup.
Namun, dengan mengambil hasil bagi oleh $\ker \phi$ dan menggunakan pemetaan $R/\ker \phi \to R'$, kita
dapat mereduksi ke kasus ketika $\phi$ injektif; dalam kasus tersebut, pernyataan ini mengikuti
dari hasil surjektivitas yang telah dibuktikan.
\end{proof}

Secara umum, terdapat \emph{banyak} pengangkatan dari suatu ideal prima tertentu.
Sebagai contoh, pertimbangkan penyertaan $\mathbb{Z} \subset \mathbb{Z}[i]$.
Ideal prima $(5) \in \spec \mathbb{Z}$ dapat diangkat menjadi $(2+i)
\in \spec \mathbb{Z}[i]$
atau $(2-i) \in \spec \mathbb{Z}[i]$.
Keduanya ideal prima yang berbeda: $\frac{2+i}{2-i} \notin \mathbb{Z}[i]$. Akan tetapi,
perhatikan bahwa setiap elemen $\mathbb{Z}$ yang habis dibagi $2+i$ secara otomatis
habis dibagi konjugatnya $2-i$, dan oleh karena itu habis dibagi hasil kali keduanya, $5$
(karena $\mathbb{Z}[i]$ merupakan UFD, sebab ia domain Euklides).

Meskipun demikian, berbagai pengangkatan tersebut tidak dapat dibandingkan.

\begin{proposition} \label{incomparableintegral}
Misalkan $\phi: R \to R'$ suatu morfisme integral. Untuk $\mathfrak{p} \in
\spec R$ yang diberikan, tidak terdapat relasi penyertaan di antara elemen-elemen $\mathfrak{q} \in \spec
R'$ yang mengangkat $\mathfrak{p}$.
\end{proposition} 
Dengan kata lain, jika $\mathfrak{q}, \mathfrak{q}' \in \spec R'$ mengangkat
$\mathfrak{p}$, maka $\mathfrak{q} \not\subset \mathfrak{q}'$.
\begin{proof} 
Kita akan memberikan pembuktian yang “licin” melalui beberapa reduksi.
Perhatikan bahwa operasi pelokalan dan pengambilan hasil bagi hanya memperkecil
$\spec$: operasi tersebut tidak “menggabungkan” ideal-ideal prima yang sebelumnya berbeda menjadi satu.
Jadi, dengan mengambil hasil bagi $R$ oleh $\mathfrak{p}$, kita boleh mengandaikan bahwa $R$ suatu
\emph{domain} dan $\mathfrak{p} = 0$.
Andaikan terdapat dua ideal prima $\mathfrak{q} \subsetneq \mathfrak{q}'$ dari $R'$
yang mengangkat $(0) \in \spec R$.
Dengan mengambil hasil bagi $R'$ oleh $\mathfrak{q}$, kita boleh mengandaikan bahwa $\mathfrak{q}  = 0$.
Kita bahkan dapat mengandaikan $R \subset R'$ dengan mengambil hasil bagi oleh kernel $\phi$.
Lema berikut dengan demikian menyelesaikan pembuktian, karena menunjukkan bahwa $\mathfrak{q}'
= 0$, suatu kontradiksi.
\end{proof} 

\begin{lemma} 
Misalkan $R \subset R'$ suatu penyertaan domain integral yang merupakan morfisme
integral. Jika $\mathfrak{q} \in \spec R'$ merupakan ideal prima tak nol, maka
$\mathfrak{q} \cap R$ tak nol.
\end{lemma} 
\begin{proof} 
Ambil $x \in \mathfrak{q}'$ yang tak nol. Berdasarkan asumsi, terdapat persamaan
\[ x^n + r_1 x^{n-1} + \dots + r_n = 0,  \quad r_i \in R, \]
yang dipenuhi oleh $x$. 
Di sini kita dapat mengandaikan $r_n \neq 0$; dengan langsung memperhatikan persamaan, diperoleh $r_n \in \mathfrak{q}' \cap R$.
Jadi, irisan ini tak nol.

\end{proof} 

\begin{corollary} 
Misalkan $R \subset R'$ suatu penyertaan domain integral sedemikian sehingga $R'$
integral di atas $R$. Jika salah satu dari $R, R'$ merupakan medan, maka yang lainnya juga.
\end{corollary} 
\begin{proof} 
Memang, $\spec R' \to \spec R$ surjektif berdasarkan \cref{lyingover}; jadi, jika $\spec
R'$ mempunyai satu elemen (yakni $R'$ merupakan medan), hal yang sama berlaku untuk $\spec R$
(yakni $R$ merupakan medan). Sebaliknya, andaikan $R$ suatu medan. 
Maka, setiap dua ideal prima dalam $\spec R'$ ditarik kembali ke elemen yang sama di $\spec
R$. Jadi, berdasarkan \cref{incomparableintegral}, tidak mungkin terdapat relasi penyertaan di antara ideal-ideal prima $\spec
R'$. Namun, $R'$ suatu domain, sehingga harus merupakan medan.
\end{proof} 


\begin{exercise} Misalkan $k$ suatu medan.
Tunjukkan bahwa $k[\mathbb{Q}_{\geq 0}]$ integral di atas gelanggang polinomial $k[T]$.
Meskipun ini merupakan perluasan yang \emph{sangat besar}, ideal prima $(T)$ hanya dapat diangkat
dengan satu cara ke $\spec k[\mathbb{Q}_{\geq 0}]$.
\end{exercise} 

\begin{exercise} 
Misalkan $A \subset B$ suatu penyertaan gelanggang di atas medan berkarakteristik
$p$. Andaikan $B^p \subset A$, sehingga $B/A$ integral dalam pengertian yang sangat kuat. 
Tunjukkan bahwa pemetaan $\spec B \to \spec A$ merupakan suatu \emph{homeomorfisme.} 
\end{exercise} 



\subsection{Teorema naik}
Misalkan $R \subset R'$ suatu penyertaan gelanggang dengan $R'$ integral di atas $R$. Kita telah melihat dalam
teorema terletak di atas (\cref{lyingover}) bahwa setiap ideal prima $\mathfrak{p} \in \spec R$ 
mempunyai ideal prima $\mathfrak{q} \in \spec R'$ yang “terletak di atas” $\mathfrak{p}$, yakni
sedemikian sehingga $R \cap \mathfrak{q} = \mathfrak{p}$.
Sekarang kita ingin menunjukkan bahwa penyertaan \emph{hingga} ideal-ideal prima dapat diangkat ke $R'$.

\begin{proposition}[Teorema naik]
Misalkan $R \subset R'$ suatu penyertaan integral gelanggang.
Andaikan $\mathfrak{p}_1 \subset \mathfrak{p}_2 \subset \dots \subset
\mathfrak{p}_n \subset R$ suatu 
rantai menaik hingga ideal-ideal prima dalam $R$.
Maka terdapat rantai menaik $\mathfrak{q}_1 \subset \mathfrak{q}_2 \subset
\dots \subset \mathfrak{q}_n$ dalam $\spec R'$ yang mengangkat rantai tersebut. 

Selain itu, $\mathfrak{q}_1$ dapat dipilih sembarang di antara ideal-ideal yang mengangkat
$\mathfrak{p}_1$.
\end{proposition} 
\begin{proof} 
Dengan induksi dan teorema terletak di atas (\cref{lyingover}), cukup ditunjukkan:
\begin{quote}
Misalkan $\mathfrak{p}_1 \subset \mathfrak{p}_2$ suatu penyertaan ideal prima dalam $\spec
R$. Misalkan $\mathfrak{q}_1 \in \spec R'$ mengangkat $\mathfrak{p}_1$. Maka terdapat
$\mathfrak{q}_2 \in \spec R'$ yang memenuhi dua syarat: mengangkat
$\mathfrak{p}_2$ dan memuat $\mathfrak{q}_1$.
\end{quote}
Untuk menunjukkan hal tersebut, kita menerapkan \cref{lyingover} pada penyertaan
$R/\mathfrak{p}_1 \hookrightarrow R'/\mathfrak{q}_1$. Terdapat suatu elemen
$\spec R'/\mathfrak{q}_1$ yang mengangkat $\mathfrak{p}_2/\mathfrak{p}_1$; elemen
$\spec R'$ yang bersesuaian dapat dipilih sebagai $\mathfrak{q}_2$.
\end{proof} 

